% results.tex
% Single Results section, no subsections, no paragraph headings.
% Flowing narrative across 8 long paragraphs:
%   (1) two-pass setup + stylized facts confirmation
%   (2) HMM-WJ jump calibration + state-resolution sweep
%   (3) single-asset 8-generator benchmark, in-sample + out-of-sample + cross-asset
%   (4) multi-asset universe, per-ticker fits, OLS calibration
%   (5) six-composer protocol + metric families
%   (6) aggregate composer results + downstream VaR
%   (7) mechanism trace + branch behavior + clipping stress
%   (8) robustness (hyperparameter grid + seed sweep)

We validated the two-stage construction in two passes. The first pass
asked whether the hybrid hidden Markov framework reproduced
single-asset stylized facts on real SPY data, in-sample and
out-of-sample, against six baseline generators. The second pass asked
whether the variance-corrected single-index composer preserved those
validated marginals when combined with a market factor across a
$424$-ticker universe, against five alternative composers, and
whether the resulting paths supported a downstream Value-at-Risk
backtest at industry-standard coverage levels. Descriptive statistics
confirmed that all three stylized facts were present in both the
in-sample ($2014$--$2024$, $T = 2{,}766$) and out-of-sample ($2025$,
$T = 249$) windows (Table~\ref{tab:descriptive_stats}): excess
kurtosis of $7.7$ (IS) and $6.9$ (OoS) with Jarque-Bera rejection
($p<0.001$), persistent volatility clustering (Ljung-Box on $|G_t|$
rejected, $p<0.001$), and near-zero linear autocorrelation in raw
returns \cite{famaEMH1970}.

With the stylized facts confirmed in both evaluation windows, we
calibrated the jump mechanism. A multi-objective grid search over
$\epsilon \in [10^{-4},\, 2.5 \times 10^{-2}]$ and
$\lambda \in [10,\, 160]$, with $200$ simulated paths per grid point,
identified the optimal values $\epsilon^{*} = 10^{-4}$ and
$\lambda^{*} = 100$ (Online Appendix~\ref{sec:supp-pseudo}). The
optimal $\lambda^{*} = 100$ was an interior point of the grid,
confirming that the mean jump duration was identifiable from data;
$\epsilon^{*} = 10^{-4}$ fell at the lower boundary of the $\epsilon$
grid, indicating that the data consistently preferred rarer
tail-entry events over the range explored. At these optimal values,
the ACF of $|G_t|$ closely tracked the observed SPY ACF across lags
$1$--$252$, confirming that the jump mechanism partially reproduced
the empirical ARCH effect. The primary analysis used $N = 100$
states, but a sensitivity study over $N \in \{30, 60, 90, 100, 150,
200\}$ confirmed robustness: KS and AD pass rates remained stable for
both HMM-NJ (without jumps) and HMM-WJ (with jumps) across this
range, with every state receiving adequate statistical support
(Online Appendix~\ref{sec:supp-sensitivity}). At $N = 350$, certain
states were never visited in the historical record, establishing a
practical upper bound on state resolution.

With the model calibrated, we benchmarked HMM-WJ against seven
alternative generators on $1{,}000$ synthetic paths of $2{,}766$
trading days each (Table~\ref{tab:model_comparison},
Figure~\ref{fig:model_comparison}). Bootstrap resampling set the
non-parametric ceiling ($100\%$ KS, $100\%$ coverage), while the
Gaussian i.i.d.\ baseline ($0\%$ KS) confirmed that thin-tailed
models were inadequate. Between these anchors, every model excelled
on one quality dimension but failed on another. The GARCH(1,1)
baseline achieved the best temporal fidelity (ACF-MAE $0.031$) but a
poor distributional fit ($5.5\%$ KS), indicating a systematic shape
mismatch despite its volatility clustering properties. The Gated
Recurrent Unit (GRU) \cite{choEtAl2014gru} closed $83\%$ of the
ACF-MAE gap but produced a $0.6\%$ KS pass rate. The hidden
semi-Markov model (HSMM) \cite{bullaBulla2006} occupied a middle
ground ($82\%$ KS) but its coarse $8$-state partition could not
capture the full growth-rate distribution, as reflected by a $38\%$
kurtosis gap and only $68.7\%$ quantile coverage. No single model
captured both quality dimensions. Both HMM variants achieved the
highest distributional fidelity among all parametric models, with
HMM-WJ's pass rates $2$--$8$ percentage points below HMM-NJ's, a
cost of the jump mechanism occasionally overweighting tail states.
The Student-$t$ ($\nu = 5$) emissions closely matched the observed
kurtosis (within $1.3\%$ for HMM-WJ), a substantial improvement over
Normal emissions which underestimated kurtosis by approximately
$29\%$. Negative skewness ($-0.75$ observed) was also well
reproduced by both HMM variants, arising from the asymmetric
state-occupancy distribution inherited from the Laplace partition
rather than any explicit skewness parameter; symmetric-innovation
models (GARCH, Gaussian, Laplace) all produced near-zero skewness by
construction. On the temporal dimension, HMM-NJ scored $0.059$ on
ACF-MAE, closing only $3\%$ of the gap between the i.i.d.\ floor
($0.060$) and the GARCH ceiling ($0.031$); HMM-WJ scored $0.052$,
closing $28\%$ of the gap and the only non-GARCH non-neural model to
reduce ACF-MAE meaningfully below the i.i.d.\ floor, because
approximately $24\%$ of HMM-WJ paths contained at least one Poisson
jump event that sustained $\mathrm{ACF}(|G_t|)$ decay across all
lags while the remaining $76\%$ behaved identically to HMM-NJ.
Out-of-sample on the held-out $249$ trading days of $2025$ the same
ranking held (Table~\ref{tab:model_comparison},
Figure~\ref{fig:statistical_validation}): HMM-NJ achieved $96.7\%$
KS and HMM-WJ achieved $94.4\%$ KS, with simulated kurtosis of $6.4$
and $6.1$ against an observed $6.9$, and HMM-WJ remained the only
model to improve meaningfully on both dimensions simultaneously.
GARCH's $W_1$ deteriorated by $67\%$ from in-sample to
out-of-sample, exposing a tail misspecification that binary pass
rates obscured. Cross-asset generalization to NVDA, JNJ, and JPM
confirmed that the jump mechanism automatically adapted its duration
$\lambda^{*}$ to each asset's clustering intensity (Online
Appendix~\ref{sec:supp-cross-asset}).

Having validated the per-asset marginals on single-asset stylized
facts in both windows, we turned to the multi-asset composition
pipeline and asked whether those validated marginals plugged into a
multi-asset construction without breaking either target. We
evaluated six composers on a fixed set of per-asset JumpHMM
marginals calibrated against the real SPY growth-rate path. This
isolated the effect of the composition rule from the effect of
market-path randomness: all composers consumed the same real SPY
history and the same per-ticker innovation draw, so metric
differences across composers reflected the composition rule alone.
The ticker universe was a $424$-asset US equity/ETF set drawn from
Polygon.io daily-close histories, filtered to tickers with a
complete in-sample trading history matching the reference length of
AAPL ($T = 2{,}767$ daily closes covering January~2014 through
December~2024); filtering on the common length preserved $424$
tickers, of which the SPY total-return growth-rate series was taken
as the market path $g_m$ and the remaining $423$ tickers as
non-market assets. For each non-market ticker $i$ we fitted a
\texttt{JumpHMM} marginal to the asset's own growth-rate history,
using the default configuration ($N = 100$ Laplace-partition states,
Student-$t$ emissions with $\nu = 5$, annualized growth rates via
\texttt{excess\_growth\_rates}, risk-free rate $r_f = 0$,
$\Delta t = 1/252$); jump parameters $(\epsilon, \lambda)$ were left
at the package defaults rather than tuned per ticker, which would
improve per-ticker marginal fit but was not required for a
comparison that tests the effect of the composition rule on top of
whatever marginal the generator produces. For each non-market
ticker we then ran ordinary least squares of the asset's
growth-rate history on $g_m$, yielding the calibrated quadruple
$(\alpha_i, \beta_i, R^2_{i,\real}, \sigma_{\eps,\real,i})$. The
universe spanned $\beta \in [0.01, 1.79]$ with median $\beta = 1.00$
and quartiles $(0.81, 1.00, 1.18)$, with $R^2_{i,\real}$ ranging
from $0.001$ to $0.933$ (Fig.~\ref{fig:r2_distribution}).

The composition rules used a \emph{paired innovation} protocol
where applicable: for each non-market ticker and each replication
$r \in \{1, \dots, R\}$ we drew a single innovation path
$\teps_i^{(r)}$ of length $T$ from the ticker's JumpHMM marginal,
and the same $\teps_i^{(r)}$ was passed to the naive,
Gaussian-residual-SIM, and hybrid composers, so that differences
among those three reflected the composition rule and not generator
stochasticity. The \emph{Naive} composer set
$g_i^{(r)} = \alpha_i + \beta_i\, g_m + \teps_i^{(r)}$ with no
variance correction. The \emph{Gaussian SIM} baseline replaced the
JumpHMM draw with
$g_i^{(r)} = \alpha_i + \beta_i\, g_m + \sigma_{\eps,\real,i}\,\xi_i^{(r)}$,
$\xi_i^{(r)} \sim \mathcal{N}(0, I_T)$ drawn fresh per replication,
recovering the calibrated $R^2$ exactly but discarding any marginal
structure carried by $\teps$. The \emph{Hybrid} composer applied
Algorithm~\ref{alg:hybrid} with floor $f = 0.10$ and threshold
$R^2_{\mathrm{preserve}} = 0.80$. Three further composers drew
their innovations from dedicated residual generators, each fit to
the real OLS residual series
$e_i(t) = G_i(t) - \alpha_i - \beta_i G_m(t)$: the
\emph{JumpHMM-on-residuals} composer fitted a fresh JumpHMM
marginal to $e_i$ via a pseudo-price inversion and composed
$g_i^{(r)} = \alpha_i + \beta_i g_m + \teps_{\mathrm{resid},i}^{(r)}$
with no variance correction; the \emph{Block bootstrap} composer
drew a Politis--Romano stationary bootstrap
replicate~\cite{politisRomano1994} of the empirical residual series
with mean block length $\sqrt{T} \approx 50$ and composed
analogously; the \emph{GARCH(1,1)-$t$} composer fitted a
conditional-variance model per ticker~\cite{bollerslev1986} and
substituted a simulated residual path, with $30$ of the $423$
tickers yielding a non-stationary fit ($\alpha + \beta \ge 1$ at
the optimizer's termination) and excluded from that composer only.
Across all composers, we skipped the optional cross-sectional rank
reorder in Algorithm~\ref{alg:hybrid} so that per-ticker metrics
reflected the marginal construction in isolation; cross-sectional
dependence under copula-based reorder is treated separately and is
outside the scope of the per-ticker evaluation reported here. For
every composed series we computed three metric families against
the ticker's real growth-rate history and the real SPY path: SIM
recovery reported $(\hat\alpha_i, \hat\beta_i, \hat R^2_i)$ from
OLS of $g_i^{(r)}$ on $g_m$ along with the absolute deviations
$|\hat\beta - \beta|$ and $|\hat R^2 - R^2_{\real}|$; marginal
fidelity was assessed by Kolmogorov--Smirnov (KS) and
Anderson--Darling (AD) two-sample $p$-values against the real
marginal, Wasserstein-$1$ distance $W_1$, excess kurtosis
$\kappa$, and the classical Hill tail-index estimator
$\hat\xi = (k-1)^{-1}\sum_{i=1}^{k-1}\log(X_{(i)}/X_{(k)})$ on
the upper $5\%$ order statistics of $|g_i^{(r)}|$, so reported
Hill values are tail indices with $\hat\xi \approx 1/\alpha$ for a
Pareto-tailed distribution with shape $\alpha$; variance
preservation was the ratio
$\mathrm{Var}(g_i^{(r)}) / \sigma_{\gen,i}^2$, using the real-data
variance as a proxy for the generator's targeted variance. Pass
rates were reported at $\alpha = 0.05$, and we reported $R = 100$
replications per ticker so that aggregate metrics were averages
over $423 \times 100 = 42{,}300$ per-ticker observations per
composer; the simulation seed was fixed at $1234$.

With the protocol established, we computed per-asset medians of
each metric and aggregated across the universe
(Table~\ref{tab:aggregate}). All composers recovered $\beta$ to
within a median absolute error of $0.018$ to $0.022$, consistent
with the identity $\mathrm{Cov}(g, g_m)/\sigma_m^2 \to
\beta_i^{\eff}$ holding at all residual scales. The differences
between composers appeared everywhere else. The Gaussian SIM
baseline recovered $R^2$ most tightly (median error $0.008$)
because it targeted $R^2$ by construction, but it failed marginal
fidelity badly: KS pass rate $0.6\%$, AD pass rate $0.5\%$, excess
kurtosis $1.4$ (real data, $\sim 13$), Hill tail index $0.22$
(versus $\sim 0.37$ on the generator; Fig.~\ref{fig:tails}). The
naive composition retained the generator marginal but inflated its
variance by $\beta_i^2 \sigma_m^2$, so its KS and AD pass rates
sat at $4.5\%$ and $2.7\%$. The hybrid composer held both
properties at once: KS pass rate $50.4\%$, AD pass rate $47.6\%$,
kurtosis $7.2$, Hill $0.36$; the $R^2$ recovery ($0.021$) came in
tighter than the naive baseline but looser than the Gaussian
baseline, reflecting that the variance-preserving branch targeted
$\sigma_{\gen,i}^2$ rather than $R^2_{i,\real}$. The three
residual-fit composers sidestepped the variance-inflation problem
entirely by fitting a generator directly to $e_i$, which by
construction has variance $\sigma_{\eps,\real,i}^2$ and no
double-counted market component; their aggregate KS pass rates
were $75.8\%$ (JumpHMM-on-residuals), $73.3\%$ (block bootstrap),
and $66.9\%$ (GARCH(1,1)-$t$), each higher than hybrid's $50.4\%$
(Fig.~\ref{fig:baseline}). The advantage sat in the body of the
distribution rather than the tails: hybrid's median excess
kurtosis ($7.17$) and Hill upper-tail index ($0.365$) matched
JumpHMM-on-residuals ($7.26$ and $0.365$) to within sampling
variability, confirming that the variance correction of
Eq.~\eqref{eq:scale} preserves the heavy-tailed character of the
full-return generator even when the rescaling is applied; block
bootstrap and GARCH, by fitting their own residual generators,
reproduced the empirical residual series but with slightly
lighter tails ($\kappa = 6.68$ and $6.04$; Hill $0.330$ and
$0.318$) than either JumpHMM-based recipe. The downstream
consequence of the body-fidelity gap was the question a
per-ticker one-day Value-at-Risk backtest was designed to answer
(Table~\ref{tab:var}, Fig.~\ref{fig:var}): at $\alpha = 0.99$ the
nominal exceedance rate is $1\%$; the naive composer
under-breached at $0.67\%$ and the Gaussian baseline
over-breached at $1.40\%$, both failing Kupiec unconditional
coverage on the majority of the universe ($45$ and $47\%$ pass
rates), while the hybrid, JumpHMM-on-residuals, block bootstrap,
and GARCH(1,1)-$t$ composers all hit coverage within one-tenth of
a percentage point of nominal ($0.95$, $0.99$, $1.09$, and
$1.09\%$ respectively), and their Kupiec pass rates clustered
between $72$ and $86\%$. Hybrid and JumpHMM-on-residuals were
statistically indistinguishable on this test: the $25$
percentage-point body-fidelity gap between them did not translate
to a Value-at-Risk accuracy difference. Hybrid's standalone
advantage is therefore scoped: it is the right rule when a
full-return per-asset generator is the given input, and it
delivers tail accuracy on par with a dedicated residual generator
without requiring a separate fitting stage.

Having established that body-fidelity did not translate into VaR
accuracy, we traced where the KS advantage over the naive
baseline came from. We plotted per-ticker KS pass rate and
per-ticker marginal variance ratio
$\mathrm{Var}(g) / \sigma_{\gen}^2$ as functions of the calibrated
$\beta$ (Fig.~\ref{fig:preservation}). The variance-ratio panel
made the mechanism explicit: the naive scatter traced the
theoretical curve $1 + \rho_i$ with
$\rho_i = \beta_i^2 \sigma_m^2 / \sigma_{\gen,i}^2$, overshooting
the generator variance by $50\%$ to $75\%$ at high $\beta$, while
hybrid and Gaussian SIM sat on the unit line by design. The
KS-pass-rate panel showed the consequence: naive fell from a pass
rate of $\sim 0.25$ at $\beta \approx 0.3$ to below $0.01$ past
$\beta \approx 0.7$, while hybrid sat between $0.3$ and $0.9$
across the full $\beta$ range, and Gaussian SIM was pinned to
zero throughout because with its Gaussian residual it did not
match the heavy-tailed real marginal at any $\beta$. The two
baselines' failure modes separated on tail metrics as well
(Fig.~\ref{fig:tails}): naive and hybrid clustered together near
the real-data kurtosis reference line
($\kappa_{\real} \approx 13$), with the difference between them
reflecting the variance inflation under naive rather than any
difference in tail shape, while Gaussian SIM collapsed to
$\kappa \approx 1$ and Hill $\approx 0.22$, pinned below the
generator regardless of $\beta$. Applying the branch-selection
rule of Algorithm~\ref{alg:hybrid} to the $423$-ticker universe
produced $421$ assignments to the \texttt{hybrid} branch, $2$ to
the \texttt{r2-preserve} branch, and no assignments to the
\texttt{hybrid-clipped} branch (Fig.~\ref{fig:branch_map},
Table~\ref{tab:branch}); the two $R^2$-preserving tickers were
QQQ ($R^2 = 0.861$) and SPYG ($R^2 = 0.933$), both ETFs that
track market-capitalization indices closely, and the
highest-$R^2$ individual stock (BLK at $R^2 = 0.645$) sat
$\sim 0.22$ below the $R^2_{\mathrm{preserve}} = 0.80$ threshold,
so the branch assignment was empirically tracker-specific on
this universe rather than a threshold artifact. The two
trackers' KS pass rate under hybrid composition was $99\%$,
consistent with the $R^2$-preserving branch recovering both the
calibrated $\beta$ and the calibrated $R^2$ by construction.
Clipping did not trigger anywhere in the universe: the market
loading ratio $\rho_i$ stayed below $1 - f = 0.90$ for every
ticker, because the JumpHMM marginals carried enough variance
that the market component never crowded out the $f = 0.10$
idiosyncratic floor. Because the $R^2$-preserve branch had only
two empirical assignments, we validated it on synthetic trackers
where the target $(\beta, R^2)$ was known by construction
(Fig.~\ref{fig:synth_tracker}): across $15$ cells of
$(\beta_{\mathrm{true}}, R^2_{\mathrm{true}}) \in
\{0.8, 1.0, 1.2\} \times \{0.80, 0.85, 0.90, 0.95, 0.99\}$, the
branch selector assigned \texttt{r2-preserve} as intended,
recovered median $\hat\beta$ to within $0.001$ of
$\beta_{\mathrm{true}}$ and median $\hat{R}^2$ to within $0.001$
of $R^2_{\mathrm{true}}$, and cleared the $\alpha = 0.05$ KS
test on at least $99\%$ of replications per cell. To check
whether the preservation advantage held across the $\beta$
distribution, we bucketed the universe by $\beta$ quartile and
recomputed the aggregate metrics within each bucket
(Table~\ref{tab:beta_bucket}): the hybrid KS pass rate declined
gently from $67\%$ in the low-$\beta$ quartile to $44\%$ in the
high-$\beta$ quartile, reflecting the caveat noted alongside
Algorithm~\ref{alg:hybrid} that as $\beta^{\eff\,2}\sigma_m^2$
grows toward $(1 - f)\sigma_{\gen}^2$ the market-driven component
of $g$ tightens the marginal toward Gaussian along its own axis;
the effect was small enough that hybrid still beat both
baselines by more than an order of magnitude in every bucket.
To probe the clipping branch, which did not trigger anywhere in
the main run, we re-ran the hybrid composer on the same universe
with the market growth-rate path rescaled by a factor
$\gamma \in \{1, 2, 3\}$, holding the calibrated marginals fixed
and drawing $R = 100$ replications per $(\mathrm{ticker},\,
\gamma)$ pair (Fig.~\ref{fig:stress}); at $\gamma = 1$ no ticker
crossed the clipping threshold and the hybrid KS pass rate
matched the headline value ($50.3\%$), while at $\gamma = 2$ and
$\gamma = 3$ the clipping branch activated on $76.4\%$ and
$95.2\%$ of tickers respectively, and the hybrid KS pass rate
rose to $71.8\%$ and $77.3\%$ above the unclipped baseline
because the clipping branch rescales $\beta^{\eff}$ to hold the
idiosyncratic floor and prevents the market component from
crowding out the marginal. The clipping branch is therefore not
a degradation under stress but a mechanism that keeps the
composed marginal near the generator target as
$\beta^{\eff\,2}\sigma_m^2$ inflates.

Finally, we asked whether the headline results were robust to the
two hyperparameter choices of Algorithm~\ref{alg:hybrid} and to
the simulation seed. We re-ran the hybrid composer on the full
$423$-ticker universe at every point of a $5 \times 5$ grid,
$R^2_{\mathrm{preserve}} \in \{0.70, 0.75, 0.80, 0.85, 0.90\}$
and $f \in \{0.05, 0.10, 0.15, 0.20, 0.30\}$
(Fig.~\ref{fig:sensitivity}); hybrid KS pass rate ranged from
$50.24\%$ to $50.38\%$ across the $25$ cells, a spread of $0.14$
percentage points, with the $R^2$-preserve assignment stable at
QQQ and SPYG for every threshold at or below $R^2 = 0.85$ and
the floor $f$ not changing any branch assignment in the grid.
The headline number was therefore invariant to the hyperparameter
values within wide tolerances, and the operating point
$(R^2_{\mathrm{preserve}} = 0.80,\, f = 0.10)$ reflects a
principled choice of the $R^2$-branch threshold and a standard
floor value rather than a data-peeked local optimum. We also
quantified the contribution of generator stochasticity to the
headline numbers by re-running the full six-composer protocol
under $8$ random seeds (Table~\ref{tab:seed_uncertainty}). The
per-composer KS pass rates held to within fractions of a
percentage point: hybrid sat at $50.40 \pm 0.18\%$ across seeds,
JumpHMM-on-residuals at $75.61 \pm 0.11\%$, block bootstrap at
$73.24 \pm 0.09\%$, GARCH(1,1)-$t$ at $66.96 \pm 0.30\%$, naive
at $4.41 \pm 0.07\%$, and Gaussian SIM at $0.60 \pm 0.04\%$. The
across-composer differences therefore exceeded the seed-induced
spread by roughly two orders of magnitude, and the seed sweep
reproduced the relative ordering of the composers without
inversions across seeds, confirming that the hybrid-vs-baseline
gap and the residual-fit-vs-hybrid body-fidelity gap were not
artifacts of the simulation seed.
